\begin{conclusion}[二元投影词迹函数的二进--Fibonacci 正规形与二面体不变性]\label{con:xi-terminal-replica-softcore-binary-wordtrace-dihedral-normal-form}
设
\[
J:=\begin{pmatrix}1&1\\[2pt]1&1\end{pmatrix},
\qquad
M:=\begin{pmatrix}1&1\\[2pt]1&0\end{pmatrix}.
\]
对任意二元词 \(\omega=\omega_1\cdots \omega_m\in\{0,1\}^m\)，定义
\[
A(\omega):=\prod_{t=1}^{m}A_{\omega_t},
\qquad
A_0:=J,\ \ A_1:=M,
\qquad
\tau(\omega):=\Tr(A(\omega)).
\]
若 \(\omega=1^m\)，则 \(\tau(\omega)=L_m\)。若 \(\omega\neq 1^m\)，则存在整数
\(
r\ge 1,\ b_1,\dots,b_r\ge 1,\ a_1,\dots,a_r\ge 0
\)，使 \(\omega\) 在循环意义下可写为
\[
\omega\sim 0^{b_1}1^{a_1}0^{b_2}1^{a_2}\cdots 0^{b_r}1^{a_r}.
\]
记 \(z:=\sum_{i=1}^{r}b_i\)（零的总数），则有显式迹公式
\[
\tau(\omega)=2^{z-r}\prod_{i=1}^{r}F_{a_i+3}.
\]
进一步，\(\tau\) 在 \(\{0,1\}^m\) 上满足循环不变性与反转不变性：
\[
\tau(\omega)=\tau(\sigma\omega),\qquad
\tau(\omega)=\tau(\omega^{\mathrm{rev}}),
\]
其中 \(\sigma\) 为任意循环移位。因而 \(\tau\) 仅依赖于三项数据：零块个数 \(r\)、零总数 \(z\)、以及一块长度多重集 \(\{a_1,\dots,a_r\}\)。
\end{conclusion}
\begin{proof}[证明要点]
令 \(u:=(1,1)^{\top}\)，则 \(J=uu^{\top}\) 且 \(J^k=2^{k-1}J\)（\(k\ge 1\)）。
当 \(\omega\neq 1^m\) 时，\(A(\omega)\) 至少含一个 \(J\)，故 \(\det A(\omega)=0\) 且 \(\mathrm{rank}(A(\omega))\le 1\)。
按所示循环分解，
\[
A(\omega)=J^{b_1}M^{a_1}\cdots J^{b_r}M^{a_r}
=2^{z-r}\,J M^{a_1}J M^{a_2}\cdots J M^{a_r}.
\]
利用 \(J=uu^{\top}\) 有
\(
J M^a J=(u^{\top}M^a u)\,J
\)，迭代得
\[
\tau(\omega)=2^{z-r}\prod_{i=1}^{r}(u^{\top}M^{a_i}u).
\]
由 Fibonacci \(Q\)-矩阵恒等式
\[
M^n=
\begin{pmatrix}
F_{n+1}&F_n\\[2pt]
F_n&F_{n-1}
\end{pmatrix},
\]
得到 \(u^{\top}M^a u=F_{a+3}\)，从而结论成立。
循环不变性由迹的循环置换不变性给出；反转不变性可由
\(\Tr(X)=\Tr(X^{\top})\) 与上式仅依赖 \(\{a_i\}\) 的对称结构推出。证毕。
\end{proof}

\begin{conclusion}[\texorpdfstring{$\Tr((T^{(q)})^m)$}{Tr((T^(q))^m)} 的全阶词迹--张量幂闭式（互补编码）]\label{con:xi-terminal-replica-softcore-trace-moment-wordtrace-tensor-closed-complementary-coding}
令 \(V_q:=\{0,1\}^q\)，定义
\[
T^{(q)}_{u,v}:=2^{-\mathbf 1_{\{u\wedge v\neq 0\}}}\in\left\{1,\frac12\right\}.
\]
则对任意 \(q\ge 2\)、\(m\ge 1\)，有严格恒等式
\[
\Tr\!\bigl((T^{(q)})^m\bigr)
=2^{-m}\sum_{\omega\in\{0,1\}^m}\tau(\omega)^q,
\]
其中 \(\tau(\omega)\) 由结论 \ref{con:xi-terminal-replica-softcore-binary-wordtrace-dihedral-normal-form} 定义。
\end{conclusion}
\begin{proof}[证明要点]
由
\(
T^{(q)}_{u,v}=\frac12+\frac12\mathbf 1_{\{u\wedge v=0\}}
\)
可得
\[
T^{(q)}=\frac12\bigl(J^{\otimes q}+M^{\otimes q}\bigr).
\]
故
\[
\bigl(J^{\otimes q}+M^{\otimes q}\bigr)^m
=\sum_{\omega\in\{0,1\}^m}\bigl(A_{\omega_1}\cdots A_{\omega_m}\bigr)^{\otimes q}
=\sum_{\omega\in\{0,1\}^m}A(\omega)^{\otimes q}.
\]
取迹并用 \(\Tr(B^{\otimes q})=(\Tr B)^q\) 得
\[
\Tr\!\bigl((T^{(q)})^m\bigr)
=2^{-m}\sum_{\omega\in\{0,1\}^m}\Tr(A(\omega))^q
=2^{-m}\sum_{\omega\in\{0,1\}^m}\tau(\omega)^q.
\]
证毕。
\end{proof}

\begin{conclusion}[例外幂和的单词替换原理与对称幂角色分离]\label{con:xi-terminal-replica-softcore-exceptional-power-sum-word-substitution-principle}
记
\[
S_m(q):=\sum_{i=0}^{q}\bigl(\lambda_i^{(q)}\bigr)^m,
\]
其中 \(\{\lambda_i^{(q)}\}_{i=0}^{q}\) 为 \(T^{(q)}\) 在 \(S_q\)-不变子空间上的例外特征值。
则对任意 \(q\ge 2\)、\(m\ge 1\)，成立
\[
S_m(q)
=2^{-m}\!\left(
\Omega_m(q)+
\sum_{\omega\in\{0,1\}^m\setminus\{1^m\}}\tau(\omega)^q
\right),
\]
其中 \(\Omega_m(q)\) 的定义见结论
\ref{con:xi-terminal-replica-softcore-lucas-pell-omega-unified}。
等价地，
\[
S_m(q)-\Tr\!\bigl((T^{(q)})^m\bigr)
=2^{-m}\bigl(\Omega_m(q)-L_m^q\bigr).
\]
\end{conclusion}
\begin{proof}[证明要点]
由结论 \ref{con:xi-terminal-replica-softcore-exceptional-spectrum-sympower-model}，
\[
T^{(q)}\big|_{\mathrm{Sym}^q}
=\frac12\bigl(\mathrm{Sym}^q(J)+\mathrm{Sym}^q(M)\bigr),
\]
故
\[
S_m(q)=2^{-m}\sum_{\omega\in\{0,1\}^m}
\Tr\!\bigl(\mathrm{Sym}^q(A(\omega))\bigr).
\]
当 \(\omega\neq 1^m\) 时，\(A(\omega)\) 含 \(J\)，故 \(\det A(\omega)=0\) 且秩为 \(1\)，特征值为 \(\{\tau(\omega),0\}\)，于是
\[
\Tr\!\bigl(\mathrm{Sym}^q(A(\omega))\bigr)=\tau(\omega)^q.
\]
当 \(\omega=1^m\) 时，\(A(\omega)=M^m\)，其贡献即
\(
\Tr(\mathrm{Sym}^q(M^m))=\Omega_m(q)
\)。
故得第一式。再与结论
\ref{con:xi-terminal-replica-softcore-trace-moment-wordtrace-tensor-closed-complementary-coding}
相减，即得第二式。证毕。
\end{proof}

\begin{corollary}[Fibonacci 原始素因子导出的高指标刚性]\label{cor:xi-terminal-replica-softcore-fibonacci-primitive-divisor-high-index-rigidity}
\leanverified{paper\_xi\_terminal\_replica\_softcore\_fibonacci\_primitive\_divisor\_high\_index\_rigidity}
令 \(\omega,\omega'\in\{0,1\}^m\setminus\{1^m\}\)，并按结论
\ref{con:xi-terminal-replica-softcore-binary-wordtrace-dihedral-normal-form}
写成
\[
\tau(\omega)=2^{z-r}\prod_{i=1}^{r}F_{a_i+3},
\qquad
\tau(\omega')=2^{z'-r'}\prod_{j=1}^{r'}F_{a'_j+3}.
\]
若 \(\tau(\omega)=\tau(\omega')\)，则两侧下标多重集在截断区间 \((12,\infty)\) 上完全一致（含重数）：
\[
\{a_i+3:\ a_i+3>12\}_{i=1}^{r}
=
\{a'_j+3:\ a'_j+3>12\}_{j=1}^{r'}.
\]
等价地，\(\{a_i:\ a_i>9\}\) 与 \(\{a'_j:\ a'_j>9\}\) 作为多重集一致。
\end{corollary}
\begin{proof}[证明要点]
令
\(
N:=\max\!\bigl(\{a_i+3\}_{i=1}^{r}\cup\{a'_j+3\}_{j=1}^{r'}\bigr)
\)。
若 \(N>12\)，由 Fibonacci 原始素因子定理（除低阶例外 \(6,12\) 外，\(F_N\) 含不整除任一更小 \(F_k\) 的素因子）可取素数 \(p\mid F_N\)，且 \(p\nmid F_k\)（\(k<N\)）。
于是等式 \(\tau(\omega)=\tau(\omega')\) 两端中，\(p\) 只能来自下标 \(N\) 的 Fibonacci 因子，故 \(F_N\) 的出现次数在两端相同。
约去共同的 \(F_N\) 因子后重复该“最大下标下降”步骤，直至所有剩余下标 \(\le 12\)。
故结论成立。证毕。
\end{proof}

\begin{proposition}[迹矩与例外幂和的二面体轨道压缩表示]\label{prop:xi-terminal-replica-softcore-dihedral-orbit-compression-trace-moments}
记 \(\mathrm{Brace}^{(2)}_m\) 为长度 \(m\) 的二元手镯集合（即 \(\{0,1\}^m\) 在循环移位与反转生成的二面体作用下的等价类），
对 \(b\in \mathrm{Brace}^{(2)}_m\) 记其轨道大小为 \(|\mathrm{Orb}(b)|\)。
则对任意 \(q\ge 2\)、\(m\ge 1\)，有
\[
\Tr\!\bigl((T^{(q)})^m\bigr)
=2^{-m}\sum_{b\in \mathrm{Brace}^{(2)}_m}
|\mathrm{Orb}(b)|\,\tau(b)^q,
\]
以及
\[
S_m(q)
=2^{-m}\!\left(
\Omega_m(q)+
\sum_{b\in \mathrm{Brace}^{(2)}_m\setminus\{[1^m]\}}
|\mathrm{Orb}(b)|\,\tau(b)^q
\right).
\]
\end{proposition}
\begin{proof}
由结论 \ref{con:xi-terminal-replica-softcore-binary-wordtrace-dihedral-normal-form}，
\(\tau\) 在二面体轨道上常值。
将结论
\ref{con:xi-terminal-replica-softcore-trace-moment-wordtrace-tensor-closed-complementary-coding}
与结论
\ref{con:xi-terminal-replica-softcore-exceptional-power-sum-word-substitution-principle}
中的词求和按二面体轨道分组，即得两式。证毕。
\end{proof}

\endinput
